% &%%%%%%%%%%%%%%%%%%%%%%%%%%%%%%%%%%%%%%%%%%%%%%%%%%%%%%%%%%%
% &%%                                                       %%
% &%%                        RESUME                         %%
% &%%                                                       %%
% &%%%%%%%%%%%%%%%%%%%%%%%%%%%%%%%%%%%%%%%%%%%%%%%%%%%%%%%%%%%

% !===========================================================

{
    \fontfamily{lmr}
    \renewcommand{\bfdefault}{bx}
    \selectfont

% ?-----------------------------------------------------------
% ? CADRE EXTERNE
% ?-----------------------------------------------------------

\makeatletter
    \newcommand\MakeBackgroundResume{
        \setlength{\@pdrLinethickness}{1pt}                                                   % Définit l'épaisseur des lignes du cadre
        \setlength{\@pdrMargeSize}{15mm}                                                      % Définit la taille des marges intérieures
        \setlength{\@pdrInnerBoxWidth}{\paperwidth-2\@pdrMargeSize-2\@pdrLinethickness}       % Calcule la largeur intérieure du cadre
        \setlength{\@pdrInnerBoxHeight}{\paperheight-2\@pdrMargeSize-2\@pdrLinethickness}     % Calcule la hauteur intérieure du cadre
        \put(\LenToUnit{\@pdrMargeSize},\LenToUnit{\@pdrMargeSize}){                          % Place le cadre aux coordonnées avec la couleur définie
            \color{abscolor}                                                                  % Applique la couleur définie pour le cadre
            \framebox(\LenToUnit{\@pdrInnerBoxWidth},\LenToUnit{\@pdrInnerBoxHeight}){}       % Dessine le cadre
        }
    }
\makeatother

% ?-----------------------------------------------------------
% ? GEOMETRIE DE LA PAGE
% ?-----------------------------------------------------------

\thispagestyle{empty}                                           % Supprime l'en-tête et le pied de page
\newgeometry{centering, vmargin=2cm, hmargin=2cm}               % Redéfinit les marges
\AddToShipoutPictureBG*{\MakeBackgroundResume}                  % Ajoute le cadre de fond

% ?-----------------------------------------------------------
% ? EN-TETE
% ?-----------------------------------------------------------

\begin{center}
    \bigskip
    {\huge \textcolor{abscolor}{\usefont{OT1}{phv}{b}{n}\selectfont DOCTORAT DE L'UNIVERSIT\'E DE TOULOUSE}}
\end{center}

\vfill

% ?-----------------------------------------------------------
% ? INFORMATIONS DE LA THESE
% ?-----------------------------------------------------------
{\parskip=0pt
{\raggedright\bfseries Délivrée par :} {\itshape ISAE - SUPAÉRO}

{\raggedright\bfseries École doctorale et spécialité :} {\itshape MEGeP -- Dynamique des fluides}

{\raggedright\bfseries Unité de Recherche :} {\itshape ISAE-ONERA EDyF Énergétique et Dynamique des Fluides}
}

{\raggedright\rule{\textwidth}{2.5pt}}\\[-2ex]
{\raggedright\rule{\textwidth}{1.2pt}}\\[-0.7cm]
\begin{center}
    {\bfseries \fontfamily{cmr} Présentée et soutenue le} \textit{\theSoutenance} {\bfseries \fontfamily{cmr} par } \textbf{\theAuteur}
\end{center}
\vspace*{-0.4cm}
{\raggedright\rule{\textwidth}{1.2pt}}\\[-1.6ex]
{\raggedright\rule{\textwidth}{2.5pt}}

\vfill

% ?-----------------------------------------------------------
% ? RESUME EN FRANÇAIS
% ?-----------------------------------------------------------

{\raggedright\bfseries\color{abscolor} Titre:} {\bfseries \theTitre}

\noindent \lipsum[1-2]

{\raggedright\bfseries\color{abscolor} Mots-clés:} LaTeX, Doctorants, DMPE, Simulation numérique, Mesure expérimentale

% ?-----------------------------------------------------------
% ? SEPARATEUR
% ?-----------------------------------------------------------

\vfill

\begin{center}
    \vspace{-10px}
    \color{abscolor}\underline{\hspace{5cm}}
\end{center}

\vfill

% ?-----------------------------------------------------------
% ? RESUME EN ANGLAIS
% ?-----------------------------------------------------------

{\raggedright\bfseries\color{abscolor} Title:} {\bfseries Votre titre de thèse en anglais}

\noindent \lipsum[1-2]

{\raggedright\bfseries\color{abscolor} Keywords:} LaTeX, PhD students, DMPE, Numerical Simulation, Experimental Measurement

\vfill

\restoregeometry%
\pagestyle{headings}

}